\documentclass[11pt,letterpaper]{article}
\usepackage[margin=0.85in]{geometry}
\usepackage{lmodern}
\usepackage{hyperref}
\hypersetup{colorlinks=false, urlcolor=black, pdfborder=0 0 0}
\setlength{\parindent}{0pt}
\setlength{\parskip}{0.45em}
\raggedright
\pagestyle{empty}
\begin{document}

\noindent
{\bfseries Ke-Cheng Hsu (Clio)}\\
Ph.D.\ Candidate, Department of Economics, The Hong Kong University of Science and Technology (HKUST)\\
(852) 6840\,9753 \quad $|$ \quad \href{mailto:khsu@connect.ust.hk}{khsu@connect.ust.hk}\\
\href{https://cliohsu1997.github.io/}{cliohsu1997.github.io} \quad $|$ \quad \href{https://github.com/cliohsu1997}{github.com/cliohsu1997}

\vspace{0.7em}
\noindent
October 26, 2026

\vspace{0.35em}
\noindent
Hiring Team\\
The Hong Kong Jockey Club --- People \& OD\\
Re: Management Trainee Programme 2027

\vspace{0.2em}
Dear Hiring Team,

I am applying for the 2027 Management Trainee Programme, with Customer and Product Development as my preferred function and Membership and Customer Excellence second. I am a Ph.D.\ candidate in Economics at HKUST, expected to graduate in June 2027, with an MPhil in Economics and a BSc in Mathematics and Economics from the same university, where I graduated first in class (CGA 3.9/4.3; Academic Achievement Medal, top 1\%). I am a Hong Kong permanent resident, a native Mandarin speaker, fluent in English, with communicative Cantonese. The Club looks for people who learn fast, solve problems, lead, and communicate. Those are the parts of my work I rely on most.

On learning and problem-solving. My research uses data to answer a decision a business faces. For an online apparel seller on Taobao from 2017 to 2023 (six categories, about 62 million RMB in revenue, return rates of about 20 to 30 percent), the decision was how lenient the return policy should be. A lenient policy, where the seller pays the buyer's return shipping fee, brings more sales but also more returns; a strict one, where the buyer pays, does the opposite. I built the data in Python and estimated which is better. The answer depends on the margin and on how much a return costs to handle: covering the fee suits a high-margin category, while in a thin-margin one, such as a basic T-shirt, letting each buyer pay for the protection improves welfare by about 50 percentage points. I then proposed a different policy --- a lower price with a partial restocking fee --- and showed it beats both, reaching about 75 percent of the best outcome (seller profit and what buyers get, together) in most categories. A back-of-envelope calculation suggests it could improve on the seller's current outcome by about 5 million RMB over the period. The same approach answers what a product team asks before a launch: how much a product with certain attributes would be used, and what it would contribute. It also shows that a policy can raise profit and still leave some customers worse off, the trade-off the Club manages when it grows wagering in a socially responsible way. I work in Python, R, and C++, and I check the data before I trust a number.

On leadership and communication. As Hall Tutor at UG Hall I, I helped residents work through conflicts. As Teaching Assistant Coordinator for the Business School in 2025--2026, I evaluated teaching assistants' mock teaching, gave feedback on what to change, and helped them improve. My own teaching is rated by students at the end of each course. Students gave me 4.53 out of 5 for lecturing the MSc Math Camp, and an average of 4.72 out of 5 across the five courses I taught as a teaching assistant. That is the score students give for how clearly the material is explained.

That is what I would like to bring to the Club's rotations and stretch projects. My CV is attached. Thank you for reading.

\vspace{0.7em}
Sincerely,\\[1.0em]
Ke-Cheng Hsu (Clio)

\end{document}
